
\sgaStableTarget{sga6.e10.prop.2-4-4}\textbf{Proposition 2.4.4.} \uline{The polynomial map} @@M0@@ (2.3.1) \uline{is constant on the cosets modulo} @@M1@@.

\emph{Proof}: If @@M2@@ is an algebraically closed extension of @@M3@@, one obtains a commutative diagram
@@M4@@
where @@M5@@ is the polynomial map on @@M7@@ corresponding to @@M6@@. Thus one may suppose that @@M8@@ is algebraically closed. The assertion then follows immediately from 2.1.3 and 2.4.3.

\sgaStableTarget{sga6.e10.rem.2-4-5}\textbf{Remark 2.4.5.} In Exp. XIII this result will be proved without using the representability of the functor @@M9@@.

\sgaStableTarget{sga6.e10.sec.3}\section*{3. Projection formulas for the associated graded groups}

\sgaStableTarget{sga6.e10.sec.3.1}\subsection*{3.1.}
Consider a noetherian scheme @@M10@@. Denote by @@M11@@ the graded group associated with the @@M12@@-filtration of @@M13@@, and by @@M14@@ the graded group associated with the filtration 1.1.1 of @@M15@@. By 1.3.2, the composition law (1.3.1)
@@M16@@
induces, on passing to the associated graded groups, composition laws
\sgaStableTarget{sga6.e10.eq.3-1-1}@@M17@@

